\documentclass{article}
\usepackage{graphicx}
\usepackage[utf8]{inputenc}
\usepackage[magyar]{babel}
\usepackage{amsmath}
\usepackage{amssymb}

\title{GBRP-PPO: Csoport-Alapú Meta-Regularizációs Keretrendszer a Jutalom-Hackelés Detektálására és Enyhítésére}

\begin{document}

\maketitle

\section{Bevezető Gondolatok}

A PPO (Proximal Policy Optimization) stabil és hatékony optimalizáló, ugyanakkor ismert jelenség, hogy exploitálható jutalomstruktúrák esetén az agent a célfeladat szemantikájától eltérő, jutalom-maximalizáló viselkedést tanulhat (reward hacking). Fontos megkülönböztetni az optimalizáció \textit{numerikus stabilitását} a tanult viselkedés \textit{szemantikai helyességétől}: reward hacking esetén a PPO tanulása lehet stabil (pl. policy loss, KL divergencia), miközben a rendszer egy hibás célstruktúrát optimalizál.

A GBRP (Group-Based Reward Protection) célja, hogy a tanulási folyamat során \textbf{csoport-alapú} és \textbf{multi-metrikus} detektálással azonosítsa a jutalom-hackelés korai jeleit, majd konfirmált esetben adaptív büntetéssel (penalty) enyhítse azt. A módszer \textbf{meta-regularizációs rétegként} működik: a PPO algoritmust nem módosítja, csak a jutalomjelet alakítja át a tanulás előtt.

\subsection*{Jelölések (implementáció-barát)}
Legyen:
\begin{itemize}
    \item $g$ a group index (policy csoport index),
    \item $k$ a policy index a group-on belül,
    \item $t$ időlépés a rollout során,
    \item $G_g$ a $g$-edik policy group, $|G|$ a group méret.
\end{itemize}

\section{Tervezési Döntések}

\subsection{Mikor kezdjük alkalmazni a penaltyt?}
\textbf{Azonnal, amikor a detektálás aktiválódik.} A grace period növeli annak esélyét, hogy az exploit rögzül.

\subsection{Visszamegyünk a history-ban?}
\textbf{Nem.} Csak forward-looking büntetés. A korábbi groupok már beépültek a policy-be.

\subsection{Mi a baseline?}
\textbf{Rolling window:} az előző $m=3$ group átlaga.

\subsection{Van recovery mechanizmus?}
\textbf{Adaptív penalty:} automatikusan gyengül, ha normalizálódnak a metrikák.

\subsection{A Self-Deception Probléma}
Ha a detektálás módosított jutalmakat használna, a rendszer a saját büntetésétől ``megnyugodna''. \textbf{Megoldás:} minden detektálás \textbf{nyers} jutalmakon és nyers metrikákon történik.

\section{Architektúra Áttekintés}

\subsection{Rendszer Felépítés}
\begin{center}
\begin{tabular}{c}
PPO $\leftrightarrow$ Environment (interakció) \\
$\downarrow$ \\
nyers jutalmak + metrikák \\
$\downarrow$ \\
GBRP (detektálás + penalty kiválasztás) \\
$\downarrow$ \\
módosított jutalmak \\
$\downarrow$ \\
PPO.update() \\
\end{tabular}
\end{center}

\textbf{Kulcs elv:} A GBRP a jutalmakat a gyűjtés és tanulás között interceptálja. A PPO standard algoritmus marad.

\subsection{Tanítási Folyamat}
\begin{enumerate}
    \item PPO rollout-ot gyűjt: states, actions, \textbf{nyers} jutalmak, entropia, policy/value loss, approx KL, stb.
    \item GBRP group-szinten kiértékeli a jelzőket (\textbf{mindig nyers} adatokból) és meghatározza a $\lambda_g$ penaltyt.
    \item PPO frissít a módosított jutalmakkal.
    \item GBRP a következő grouphoz frissíti a rolling baseline-okat (szintén \textbf{nyers} adatokból).
\end{enumerate}

\section{Algoritmus Specifikáció}

\subsection{Group Struktúra}
\begin{itemize}
    \item Group méret: $|G| = 5$ policy
    \item Detektálás: minden group lezárása után (group-level döntés)
    \item Penalty alkalmazás: forward-looking (a következő group rewardjára/rolloutjaira)
\end{itemize}

\subsection{Baseline számítás (nyers)}
Legyen:
\begin{equation}
\bar{R}_g^{\text{raw}} = \frac{1}{|G|}\sum_{k\in G_g} R_{g,k}^{\text{raw}},
\quad
\bar{L}_g = \frac{1}{|G|}\sum_{k\in G_g} L_{g,k}
\end{equation}

Rolling baseline (előző $m=3$ group):
\begin{equation}
\mu_R^{(g-1)} = \frac{1}{m}\sum_{i=g-m}^{g-1}\bar{R}_i^{\text{raw}},
\quad
\mu_L^{(g-1)} = \frac{1}{m}\sum_{i=g-m}^{g-1}\bar{L}_i
\end{equation}

\subsection{Detektálási Mechanizmus}

\textbf{KRITIKUS: Minden detektálás a NYERS jutalmakra és nyers metrikákra épül!}

\textbf{Kétszintű detektálás:} Level 1 (korai/prediktív) + Level 2 (konfirmált).

\subsubsection{Level 1 -- Preventív Detektálás (Korai Jelek / Előrejelzés)}
Cél: a reward hacking \textit{kialakulását} jelezni, mielőtt rögzülne.

\paragraph{(L1) Action Imbalance}
\begin{equation}
a_{\text{imb}}^{(g)} = \max_i \left(\frac{\#(action_i)}{\#(all\ actions)}\right)
\end{equation}
Küszöb: $a_{\text{imb}}^{(g)} > \tau_a$ (pl. $\tau_a=0.6$)

\paragraph{(L1) Entropy Drop (hirtelen csökkenés)}
\begin{equation}
e_{\text{drop}}^{(g)} = \frac{H_{g-1} - H_g}{|H_{g-1}| + \epsilon}
\end{equation}
Küszöb: $e_{\text{drop}}^{(g)} > \tau_e$ (pl. $\tau_e=0.15$)

\paragraph{(L1) Reward Growth (gyors növekedés nyers rewardban)}
\begin{equation}
r_{\text{grad}}^{(g)} = \frac{\bar{R}_g^{\text{raw}} - \mu_R^{(g-1)}}{|\mu_R^{(g-1)}|+\epsilon}
\end{equation}
Küszöb: $r_{\text{grad}}^{(g)} > \tau_r$ (pl. $\tau_r=0.2$)

\paragraph{(L1) Value Loss Instability (kritikus jel)}
Legyen $L_v^{(g)}$ a value loss idősor a groupban.
\begin{equation}
v_{\text{inst}}^{(g)} = \frac{\mathrm{Var}(L_v^{(g)})}{\mathrm{Var}(L_v^{(g-1)})+\epsilon}
\end{equation}
Küszöb: $v_{\text{inst}}^{(g)} > \tau_v$ (pl. $\tau_v=1.5$)

\paragraph{Level 1 aktiválás (több jel egyszerre)}
Level 1 akkor aktív, ha legalább $n_1$ korai jel teljesül:
\begin{equation}
S_1^{(g)} =
\mathbb{1}[a_{\text{imb}}^{(g)}>\tau_a] +
\mathbb{1}[e_{\text{drop}}^{(g)}>\tau_e] +
\mathbb{1}[r_{\text{grad}}^{(g)}>\tau_r] +
\mathbb{1}[v_{\text{inst}}^{(g)}>\tau_v]
\end{equation}

\begin{equation}
\text{Level1\_active}^{(g)} = \mathbb{1}[S_1^{(g)} \ge n_1], \quad n_1=1\ \text{(konzervatív)}\ \text{vagy}\ 2\ \text{(szigorúbb)}
\end{equation}

\subsubsection{Level 2 -- Konfirmált Exploit Detektálás (multi-metrikus)}
Cél: megerősíteni, hogy a rendszer \textit{stabilan optimalizál egy rossz célt}.

\paragraph{(L2) Reward Spike Ratio}
\begin{equation}
r_{\text{spike}}^{(g)} = \frac{\bar{R}_g^{\text{raw}}}{\mu_R^{(g-1)}+\epsilon}
\end{equation}
Küszöb: $r_{\text{spike}}^{(g)} > \tau_{\text{spike}}$ (pl. $1.5$)

\paragraph{(L2) Episode Length--Reward Coupling (empirikusan erős jel)}
\begin{equation}
E_{\text{len}}^{(g)} = \mathbb{1}\left[\bar{L}_g > \mu_L^{(g-1)}\ \wedge\ \bar{R}_g^{\text{raw}} > \mu_R^{(g-1)}\right]
\end{equation}

\paragraph{(L2) Explained Variance Stagnation}
Legyen:
\begin{equation}
EV_g = 1 - \frac{\mathrm{Var}(R_g - \hat{V}_g)}{\mathrm{Var}(R_g)+\epsilon}
\end{equation}
Stagnáció jel:
\begin{equation}
E_{\text{EV}}^{(g)} = \mathbb{1}[\alpha < EV_g < \beta],\quad \alpha=0.3,\ \beta=0.7
\end{equation}

\paragraph{(L2) Entropy Oscillation (nem csak csökkenés!) }
Késői exploit-fázisban az entrópia nem feltétlen csökken tovább, hanem oszcillál:
\begin{equation}
\sigma_H^{(g)} = \mathrm{Std}(H^{(g)}),\quad
E_{\text{ent}}^{(g)}=\mathbb{1}[\sigma_H^{(g)}>\tau_H]
\end{equation}

\paragraph{(L2) Value Loss Instability (megerősítésként)}
\begin{equation}
E_{\text{v}}^{(g)} = \mathbb{1}[v_{\text{inst}}^{(g)}>\tau_v]
\end{equation}

\paragraph{(Megjegyzés) Policy loss nem önálló jel}
A policy gradient loss konvergenciája 0-hoz nem tekinthető exploit-indikátornak:
\begin{equation}
\lim_{g\to\infty}\mathcal{L}_{\text{policy}}^{(g)} \to 0 \;\nRightarrow\; \text{helyes céloptimalizálás}
\end{equation}
Reward hacking esetén a PPO lehet ``elégedett'' (nincs clipping probléma), miközben hibás jutalomszerkezetet optimalizál.

\paragraph{Level 2 aktiválás (több jel egyszerre)}
Legyen:
\begin{equation}
S_2^{(g)} =
\mathbb{1}[r_{\text{spike}}^{(g)}>\tau_{\text{spike}}] +
E_{\text{len}}^{(g)} +
E_{\text{EV}}^{(g)} +
E_{\text{ent}}^{(g)} +
E_{\text{v}}^{(g)}
\end{equation}

\begin{equation}
\text{Level2\_active}^{(g)}=\mathbb{1}[S_2^{(g)}\ge n_2],\quad n_2=2\ \text{(alap)}\ \text{vagy}\ 3\ \text{(szigorú)}
\end{equation}

\section{Adaptív Penalty (hierarchikus, több jel alapján)}

A penalty együttható $\lambda_g\in[0,1]$ a detektálás alapján:

\begin{equation}
\lambda_g =
\begin{cases}
0.4 & \text{ha Level2\_active}^{(g)}=1 \\
0.7 & \text{ha Level1\_active}^{(g)}=1 \ \wedge\ r_{\text{grad}}^{(g)} > 0.25 \\
0.9 & \text{ha Level1\_active}^{(g)}=1 \\
1.0 & \text{egyébként (recovery)}
\end{cases}
\end{equation}

Módosított jutalom (forward-looking):
\begin{equation}
r_{g,k}^{\text{mod}} = \lambda_g \cdot r_{g,k}^{\text{raw}}
\end{equation}

\textbf{Tulajdonságok:}
\begin{itemize}
    \item Csak forward-looking (nincs retrospektív módosítás)
    \item Detektálás mindig nyers adatokból (anti self-deception)
    \item Aktiválás több jel együttes teljesülése esetén
\end{itemize}

\section{Monitorozás és Metrikák}

\subsection{TensorBoard Logging}

\subsubsection{PPO Metrikák (Standard)}
\begin{itemize}
    \item \texttt{ppo/policy\_loss}
    \item \texttt{ppo/value\_loss}
    \item \texttt{ppo/entropy}
    \item \texttt{ppo/approx\_kl}
    \item \texttt{rollout/ep\_len\_mean}
    \item \texttt{rollout/ep\_rew\_mean}
\end{itemize}

\subsubsection{GBRP Metrikák (Új)}
\begin{itemize}
    \item \texttt{gbrp/level1\_active}
    \item \texttt{gbrp/level2\_active}
    \item \texttt{gbrp/S1\_count}
    \item \texttt{gbrp/S2\_count}
    \item \texttt{gbrp/current\_lambda}
    \item \texttt{gbrp/reward\_grad}
    \item \texttt{gbrp/reward\_spike}
    \item \texttt{gbrp/episode\_len\_reward\_coupling}
    \item \texttt{gbrp/value\_loss\_instability}
    \item \texttt{gbrp/explained\_variance}
    \item \texttt{gbrp/entropy\_oscillation}
\end{itemize}

\subsubsection{Duális jutalom követés}
\begin{itemize}
    \item \texttt{rewards/raw}
    \item \texttt{rewards/modified}
\end{itemize}

\section{Kulcsfontosságú Technikai Insights}

\subsection{Multi-metrikus robusztusság}
Egyetlen metrikára épülő detektálás sebezhető. A GBRP több jel együttes teljesülését követeli meg ($S_1\ge n_1$, $S_2\ge n_2$), így csökkenti a fals pozitívokat és a gaming lehetőségét.

\subsection{Stabil optimalizáció vs. rossz cél}
Reward hacking esetén a PPO stabil lehet (policy loss, KL), miközben a value loss instabil, az epizódhossz és a jutalom együtt nő, és a critic nem modellezi stabilan a returnt. A GBRP célja ezen stabil, de rossz céloptimalizálás detektálása és korrekciója.

\section{Következtetés}

A GBRP-PPO egy csoport-alapú, multi-metrikus meta-regularizációs keretrendszer, amely a reward hacking korai jeleit (Level 1) és konfirmált mintázatait (Level 2) külön kezeli. A büntetés csak több jel együttes teljesülése esetén aktiválódik, és a detektálás mindig nyers adatokból történik, megelőzve az öncsalást. A képletek közvetlenül implementálhatók callback-alapú PPO pipeline-ban.

\end{document}
